\documentclass[a4paper,11pt]{article}
        \usepackage[top=15mm,bottom=18mm,left=16mm,right=16mm]{geometry}
        \usepackage{fontspec}
        \usepackage{xeCJK}
        \setmainfont{Times New Roman}
        \setCJKmainfont{Yu Gothic}
        \setmonofont{Consolas}
        \setCJKmonofont{Yu Gothic}
        \usepackage{booktabs}
        \usepackage{longtable}
        \usepackage{tabularx}
        \usepackage{array}
        \usepackage{enumitem}
        \usepackage{hyperref}
        \usepackage{listings}
        \usepackage{xcolor}
        \usepackage{amsmath}
        \usepackage{amssymb}
        \hypersetup{
          colorlinks=true,
          linkcolor=blue,
          urlcolor=blue,
          pdftitle={sim ROS2 launch 入口},
          pdfauthor={Codex}
        }
        \lstset{
          basicstyle=\ttfamily\small,
          breaklines=true,
          columns=fullflexible,
          keepspaces=true,
          frame=single,
          framerule=0.2pt,
          rulecolor=\color{black},
          showstringspaces=false
        }
        \setlength{\parskip}{0.42em}
        \setlength{\parindent}{1em}
        \setlength{\emergencystretch}{3em}
        \renewcommand{\arraystretch}{1.15}
        \title{05. sim ROS2 launch 入口}
        \author{solar\_ws0129-main}
        \date{2026-07-29}
        \begin{document}
        \maketitle
        \tableofcontents
        \clearpage

        \section{対象}
        \begin{tabularx}{\linewidth}{>{\raggedright\arraybackslash}p{0.18\linewidth} X}
        \toprule
        項目 & 内容 \\
        \midrule
        ファイル & \path|launch/solarcar_sim.launch.py| \\
        ソースSHA-256 & \path|b674b3680c605030b54de4c2667bfa8ca8e409fbce6f6a970953bfa63f954da2| \\
        種別 & ROS 2 launch Python \\
        区分 & launch \\
        役割 & GPS 模擬、車体状態模擬、MPC、dashboard を立ち上げる simulation launch。 \\
        起動文脈 & sim モード起動時に ros2 launch される。 \\
        \bottomrule
        \end{tabularx}

        \section{このファイルを読む理由}
        GPS 模擬、車体状態模擬、MPC、dashboard を立ち上げる simulation launch。

        \begin{itemize}[leftmargin=1.6em]
        \item route/forecast/maps の存在確認を先に行う。
\item sim 用の common\_mpc パラメータ束を構成する。
\item planner speed が GPS/state 側へ戻る閉ループを作る。
        \end{itemize}

        \section{呼び出し関係}
        \subsection{どこから呼ばれるか}
        \begin{itemize}[leftmargin=1.6em]
        \item \path|scripts/solar_control.sh|
\item \path|SolarSim.ps1|
        \end{itemize}

        \subsection{次に読むと理解が進むファイル}
        \begin{itemize}[leftmargin=1.6em]
        \item \path|mpc_solarcar/gps_sim_node.py|
\item \path|mpc_solarcar/solar_state_node.py|
\item \path|mpc_solarcar/mpc_node.py|
        \end{itemize}

        \subsection{同一リポジトリ内の主要依存}
        \begin{itemize}[leftmargin=1.6em]
        \item \path|mpc_solarcar/solar_profile.py|
        \end{itemize}

        \section{ソース冒頭の把握}
        モジュール docstring または先頭意図の要約:

        モジュール docstring は明示されていない。

        \subsection{代表コード断片}
        \begin{lstlisting}[language=Python]
"s0_km": float(sim_cfg.get("start_s_km", 0.0)),
    }

    return [
        Node(
            package="mpc_solarcar",
            executable="gps_sim_node",
            name="gps_sim_node",
            parameters=[
                {
                    "route_csv": route_waypoints_csv,
                    "dt": float(sim_cfg.get("gps_rate_hz", 1.0)),
                    "init_speed_kmh": float(sim_cfg.get("gps_init_speed_kmh", sim_cfg.get("v0_kmh", 45.0))),
                }
            ],
\end{lstlisting}


        \section{主要構造の読み取り}
        主要関数は generate\_launch\_description。 launch から起動する Node action は 4 件。

        \subsection{主要クラス・関数}
            \begin{longtable}{p{0.07\linewidth} p{0.16\linewidth} p{0.25\linewidth} p{0.40\linewidth}}
            \toprule
            行 & 種別 & 名前 & 補足 \\
            \midrule
            17 & function & \path|_cfg_path| & args: profile\_path, raw \\
21 & function & \path|_setup| & args: context \\
146 & function & \path|generate_launch_description| &  \\
            \bottomrule
            \end{longtable}

        \subsection{ファイルを上から読んだときの定義順}
            \begin{longtable}{p{0.07\linewidth} p{0.84\linewidth}}
            \toprule
            行 & 内容 \\
            \midrule
            17 & 関数 \_cfg\_path を定義する。 \\
21 & 関数 \_setup を定義する。 \\
146 & 関数 generate\_launch\_description を定義する。 \\
            \bottomrule
            \end{longtable}

        \subsection{import 群の詳細}
            \begin{longtable}{p{0.07\linewidth} p{0.33\linewidth} p{0.50\linewidth}}
            \toprule
            行 & import 文 & このファイルで読む理由 \\
            \midrule
            1 & \path|from __future__ import annotations| & 型ヒントの遅延評価を有効にし、自己参照型や forward reference を安全に扱うため。 このファイル内での使用位置は少ないか、間接利用である。 \\
3 & \path|from launch import LaunchDescription| & launch が実行すべき action 群をまとめるため。 このファイル内での主な使用位置は L147。 \\
4 & \path|from launch.actions import DeclareLaunchArgument, OpaqueFunction| & DeclareLaunchArgument や OpaqueFunction など launch action を使うため。 このファイル内での主な使用位置は L149, L150。 \\
5 & \path|from launch.substitutions import LaunchConfiguration| & launch 引数の実行時値を参照するため。 このファイル内での主な使用位置は L22。 \\
6 & \path|from launch_ros.actions import Node| & ROS 2 ノード起動 action を launch から記述するため。 このファイル内での主な使用位置は L83, L95, L126, L132。 \\
8 & \path|from mpc_solarcar.solar_profile import get_path, get_section, load_profile, require_csv_data_rows, resolve_relative_path| & profile YAML 読込と検証 から get\_path, get\_section, load\_profile, require\_csv\_data\_rows, resolve\_relative\_path を読み込み、このファイルの内部処理を分担させるため。 実体ファイルは mpc\_solarcar/solar\_profile.py。 このファイル内での主な使用位置は L18, L23, L24, L25, L27, L28, L29, L30, ...。 \\
            \bottomrule
            \end{longtable}

        \section{このファイルを読む前に必要な基礎知識}
        次の章は、構文やROS用語を既知と仮定しないための説明である。

        \subsection{プログラム、プロセス、メモリ、オブジェクトを区別する}
ソースファイルはディスク上の文字列であり、それ自体は走っていない。Python実行可能プログラムがソースを読み、OSがその実行を一つのプロセスとして管理する。プロセスには仮想メモリ、開いているファイル、スレッド、終了コードなどが対応する。


\begin{lstlisting}
ソースファイル -> Pythonインタプリタが読み込む -> OS上のプロセス -> Pythonオブジェクトをメモリに生成 -> 関数やcallbackを実行
\end{lstlisting}


「メモリ上に生成する」とは、実行中プロセスが使う記憶領域に、その値の型、属性、参照関係を表す実体を用意することである。変数は実体そのものというより、そのオブジェクトを指す名前として理解するとPythonの代入が読みやすい。


\begin{lstlisting}[language=python]
node = MPCNode()
alias = node
# nodeとaliasは同じオブジェクトを参照する。
\end{lstlisting}


一つのプロセスには複数スレッドを持たせられる。スレッドは同じプロセスのメモリを共有するため、受信callbackと最適化callbackが同じself.zを書き換える場合は実行順と排他を検討する必要がある。

\paragraph{根拠資料}
\begin{itemize}[leftmargin=1.6em]
\item \href{https://docs.python.org/3/tutorial/classes.html}{Python公式チュートリアル: Classes}
\end{itemize}

\subsection{名前、代入、型、参照}
Pythonの代入文は、右辺を先に評価し、その結果のオブジェクトへ左辺の名前を結び付ける。`x = f()`では、まずfを呼び、その戻り値が得られてからxが更新される。


`float(...)`、`int(...)`、`bool(...)`は型名を呼び出して値を変換する。外部CSV、YAML、ROS parameterから来た値は型が期待どおりとは限らないため、このリポジトリでは明示変換が多い。


`None`は値が無いことを示す単一のオブジェクト、`math.nan`は浮動小数点値ではあるが有効な数値ではないことを示す。両者は用途が異なるため、`is None`と`math.isfinite`を使い分ける。

\paragraph{根拠資料}
\begin{itemize}[leftmargin=1.6em]
\item \href{https://docs.python.org/3/tutorial/classes.html}{Python公式チュートリアル: Classes}
\end{itemize}

\subsection{関数、引数、戻り値、スコープ}
`def`は関数本体をその場で実行する文ではない。関数オブジェクトを作り、その名前を現在の名前空間へ登録する。`f`は関数そのもの、`f()`はその関数を呼んだ結果である。


仮引数は関数定義側の受け取り名、実引数は呼出側から渡す値である。位置引数、キーワード引数、既定値付き引数、`*args`、`**kwargs`は渡し方が異なる。


\begin{lstlisting}[language=python]
def clip_speed(value, minimum=0.0, maximum=110.0):
    return min(maximum, max(minimum, value))

v = clip_speed(120.0, maximum=90.0)
\end{lstlisting}


関数内で作った通常の名前はローカルスコープに属する。メソッドが`self.z`を更新すればオブジェクトの状態が残るが、単に`z`を更新しただけならその呼出しのローカル値である。

\paragraph{根拠資料}
\begin{itemize}[leftmargin=1.6em]
\item \href{https://docs.python.org/3/tutorial/classes.html}{Python公式チュートリアル: Classes}
\end{itemize}

\subsection{module、package、import、console\_scripts}
Pythonファイルはmoduleとして読み込める。複数moduleをディレクトリにまとめたものがpackageである。`from .model import SolarCarModel`の先頭の点は、現在と同じpackage内のmodel moduleを指す相対importである。


import時にはファイルのトップレベル文が上から一度実行される。`def`や`class`は関数・クラスを登録するが、その本体の通常処理は呼び出すまで走らない。


setuptoolsの`console\_scripts`は、端末で使う実行可能名と`package.module:function`を対応付けるインストール時メタデータである。生成される実行用ラッパーはmoduleをimportし、指定関数を呼び、戻り値を終了コードとして扱う小さな入口である。


\begin{lstlisting}
ROS launchのexecutable名 -> install済みconsole script -> mpc_solarcar.mpc_nodeをimport -> main()を呼ぶ
\end{lstlisting}

\paragraph{根拠資料}
\begin{itemize}[leftmargin=1.6em]
\item \href{https://setuptools.pypa.io/en/latest/userguide/entry_point.html}{setuptools公式: Entry Points / Console Scripts}
\item \href{https://docs.python.org/3/tutorial/classes.html}{Python公式チュートリアル: Classes}
\end{itemize}

\subsection{例外、try/finally、with、資源解放}
例外は通常の戻り値とは別経路で異常を呼出元へ伝える。`try/except`は想定した異常を処理し、`finally`は成功・失敗にかかわらず後始末を行う。


`with open(...) as f:`はcontext managerを使い、ブロックを出るとファイルを閉じる。CSVやログの破損を避けるため、開いた資源の所有者と閉じる場所を明確にする。


`except Exception: pass`はノードを止めない利点がある一方、入力異常を隠して原因追跡を難しくする。安全に関係する値では、少なくとも頻度制限付き警告、異常カウンタ、fallback状態のpublishを検討する。



\subsection{launch Action、Node Action、実行可能名、remapping}
`launch\_ros.actions.Node(...)`はrclpyのNode基底クラスではなく、指定したpackageのexecutableをプロセスとして起動するlaunch Actionである。


`DeclareLaunchArgument`はlaunch実行時に受け取る入力欄を宣言し、`LaunchConfiguration`はその値を後で解決するsubstitutionを表す。`perform(context)`は実行時contextから確定文字列を取り出す。


launchの`name`はNode名override、`parameters`は起動時parameter、`remappings`はNodeやtopicの既定名を別名へ対応付ける。launchはPythonファイルからmainという名前を推測せず、executableとしてインストールされたconsole scriptを起動する。

\paragraph{根拠資料}
\begin{itemize}[leftmargin=1.6em]
\item \href{https://docs.ros.org/en/humble/Tutorials/Beginner-CLI-Tools/Understanding-ROS2-Nodes/Understanding-ROS2-Nodes.html}{ROS 2 Humble公式: Understanding nodes}
\item \href{https://setuptools.pypa.io/en/latest/userguide/entry_point.html}{setuptools公式: Entry Points / Console Scripts}
\end{itemize}

\subsection{天候、route、補間、時刻、単位}
予報は時刻だけの系列か、時刻とroute距離の2次元gridになり得る。現在時刻tと距離sがgrid点の間なら、周囲の値から補間して日射、温度、風を求める。


UTC、現地timezone、naive datetimeを混ぜると数時間のずれが発生する。入力でtimezoneを確定し、内部比較はUTC、表示は現地時刻という役割分担が安全である。


route距離のkm、速度のkm/hとm/s、時間の秒、energyのWhとJを跨ぐため、変換係数3.6、3600、1000の意味を各式で確認する。



\subsection{制御、状態、入力、モデル予測制御MPC}
制御対象の内部を表す状態をx、操作入力をu、外乱・予報をwとすると、離散モデルは`x[k+1] = f(x[k], u[k], w[k])`と書ける。ソーラーカーではSoC、電池温度、距離、速度などが状態候補、速度目標や駆動トルクが入力候補になる。


\[
\min_{u_0,\ldots,u_{N-1}} \sum_{k=0}^{N-1}\ell(x_k,u_k,w_k)+V_f(x_N)
\]

MPCは現在状態からNステップ先まで予測し、目的関数と制約を満たす入力系列を求める。ただし実際に適用するのは通常先頭入力だけで、次回は新しい実測状態から再び解く。これがreceding horizonである。


予測モデル、目的関数、制約、ホライズン、solver、初期値のどれかが変わると答えも変わる。「MPCを使う」だけでは仕様は決まらず、これらを単位付きで追う必要がある。

\paragraph{根拠資料}
\begin{itemize}[leftmargin=1.6em]
\item \href{https://docs.scipy.org/doc/scipy/reference/generated/scipy.optimize.minimize.html}{SciPy公式: scipy.optimize.minimize}
\end{itemize}

\subsection{ROS graph、Node、topic、service、action、parameter}
ROS graphは、実行中Nodeと、それらが持つpublisher、subscription、service、actionなどの接続関係である。Pythonクラス、プロセス、ROS graph上のNode名は関連するが同一物ではない。


topicは継続データ向けの非同期一方向publish/subscribe、serviceは短いrequest/response、actionは時間のかかる目標へfeedback、cancel、resultを持たせる。parameterはNodeごとの設定値である。


publisherが送った時点とsubscription callbackが実行される時点は同じとは限らない。通信遅延、QoS queue、executor待ちを挟むため、センサ値にはtimestampとfreshness判定が必要になる。

\paragraph{根拠資料}
\begin{itemize}[leftmargin=1.6em]
\item \href{https://docs.ros.org/en/humble/Tutorials/Beginner-CLI-Tools/Understanding-ROS2-Nodes/Understanding-ROS2-Nodes.html}{ROS 2 Humble公式: Understanding nodes}
\item \href{https://docs.ros.org/en/humble/Concepts/Basic/Interfaces-Topics-Services-Actions.html}{ROS 2 Humble公式: Topics, Services, Actions}
\item \href{https://docs.ros.org/en/humble/Concepts/Basic/About-Parameters.html}{ROS 2 Humble公式: Parameters}
\end{itemize}

\subsection{rqt\_graph、ros2 CLI、rosbag2をいつ使うか}
rqt\_graphは接続関係を見る道具であり、数値の正しさや更新周期までは保証しない。起動直後、topic名変更後、publisherが複数存在する疑いがある時に使う。


\begin{lstlisting}[language=bash]
ros2 node list
ros2 node info /mpc_node
ros2 topic list -t
ros2 topic info -v /planner/speed_cmd
ros2 topic hz /planner/speed_cmd
ros2 topic echo /planner/status
rqt_graph
\end{lstlisting}


rosbag2はtopicメッセージを時系列のまま記録・再生する。通信不具合、freshness、再計画trigger、実車とSILSの差を再現可能にするため、本番前試験では制御入力だけでなく原因となる全telemetry、status、parameter情報を記録する。


\begin{lstlisting}[language=bash]
ros2 bag record -o outputs/bags/preflight \
  /vehicle/s_km /vehicle/speed_kmh /vehicle/batt_soc \
  /vehicle/batt_temp_c /vehicle/batt_current_a /vehicle/batt_voltage_v \
  /planner/upper_speed_cmd /planner/speed_cmd /planner/status

ros2 bag info outputs/bags/preflight
ros2 bag play outputs/bags/preflight --clock
\end{lstlisting}


bag再生時はQoS互換性、simulation time、外部publisherとの二重入力に注意する。実車Nodeを同時に動かす場合はnamespaceまたはremappingで入力源を明確に分離する。

\paragraph{根拠資料}
\begin{itemize}[leftmargin=1.6em]
\item \href{https://docs.ros.org/en/ros2_packages/humble/api/rqt_graph/index.html}{ROS 2 Humble公式: rqt\_graph}
\item \href{https://docs.ros.org/en/humble/Tutorials/Beginner-CLI-Tools/Recording-And-Playing-Back-Data/Recording-And-Playing-Back-Data.html}{ROS 2 Humble公式: Recording and playing back data}
\item \href{https://docs.ros.org/en/humble/Tutorials/Beginner-CLI-Tools/Understanding-ROS2-Nodes/Understanding-ROS2-Nodes.html}{ROS 2 Humble公式: Understanding nodes}
\end{itemize}

\subsection{単体試験、SILS、replay、観測可能性}
単体試験は関数やモデル項の局所契約、SILSは複数Nodeと時間進行を含む閉ループ、historical replayは実測入力への再現性、本番preflightは実機通信と停止動作を確認する。目的が異なるため一つで代用しない。


再現可能な診断には、入力、出力、内部状態、solver成否、候補数、計算時間、fallback、source revision、profile hashを同じrun IDで記録する。


非同期不具合は最終値だけでは追えない。topic周期、message age、callback所要時間、queue遅延、Node生存、publisher数を同時に観測する。

\paragraph{根拠資料}
\begin{itemize}[leftmargin=1.6em]
\item \href{https://docs.ros.org/en/ros2_packages/humble/api/rqt_graph/index.html}{ROS 2 Humble公式: rqt\_graph}
\item \href{https://docs.ros.org/en/humble/Tutorials/Beginner-CLI-Tools/Recording-And-Playing-Back-Data/Recording-And-Playing-Back-Data.html}{ROS 2 Humble公式: Recording and playing back data}
\item \href{https://docs.ros.org/en/rolling/Concepts/Intermediate/About-Executors.html}{ROS 2公式: Executors}
\end{itemize}



        \section{関数・クラスを上から順に解説}
        \subsection{L17 関数 \_cfg\_path}
            \begin{itemize}[leftmargin=1.6em]
              \item 行範囲: L17--L18
              \item 所属: module直下
              \item 定義: \path|_cfg_path(profile_path, raw: str) -> str|
              \item docstring: 明示されていない。
              \item このブロックが直接呼ぶ主な関数・メソッド: resolve\_relative\_path, str, strip
              \item 戻り値の要点: resolve\_relative\_path(profile\_path.parent, raw) if str(raw or '').strip() else ''
              \item 主なローカル代入名: 特になし。
              \item 読み取る主なインスタンス属性: 特になし。
              \item 更新する主なインスタンス属性: 特になし。
              \item 明示的に送出する例外: 明示的raiseはない。
              \item 制御構造の規模: 条件分岐 0、ループ 0、try 0
            \end{itemize}
            \paragraph{この定義を読むためのPython構文}
            \begin{itemize}[leftmargin=1.6em]
            \item `->`以後は戻り値の型注釈であり、実行時の値を自動変換する命令ではない。
            \end{itemize}
            \paragraph{上から順の処理}
            \begin{enumerate}[leftmargin=1.8em]
            \item resolve\_relative\_path(profile\_path.parent, raw) if str(raw or '').strip() else '' を返す。
            \end{enumerate}
            \paragraph{代表コード断片}
            \begin{lstlisting}[language=Python]
def _cfg_path(profile_path, raw: str) -> str:
    return resolve_relative_path(profile_path.parent, raw) if str(raw or "").strip() else ""
\end{lstlisting}

\subsection{L21 関数 \_setup}
            \begin{itemize}[leftmargin=1.6em]
              \item 行範囲: L21--L143
              \item 所属: module直下
              \item 定義: \path|_setup(context)|
              \item docstring: 明示されていない。
              \item このブロックが直接呼ぶ主な関数・メソッド: LaunchConfiguration, Node, float, get, get\_path, get\_section, int, items, load\_profile, perform, require\_csv\_data\_rows, str
              \item 戻り値の要点: [Node(package='mpc\_solarcar', executable='gps\_sim\_node', name='gps\_sim\_node', parameters=[\{'route\_csv': route\_waypoints\_csv, 'dt': float(sim\_cfg.get('gps\_rate\_hz', 1.0)), 'init\_speed\_kmh': float(sim\_cfg.get('gps\_init\_speed\_kmh', sim\_cfg.get('v0\_kmh', 45.0)))\}]), Node(package='mpc\_solarcar', executable='solar\_state\_node', name='solar\_state\_node', parameters=[\{'profile\_yaml': str(profile\_path), 'forecast\_csv': forecast\_csv, 'route\_profile\_csv': route\_profile\_csv, 'params\_yaml': str(profile\_path), 'publish\_rate\_hz': float(sim\_cfg.get('gps\_rate\_hz', 1.0)), 'init\_speed\_kmh': float(sim\_cfg.get('v0\_kmh', sim\_cfg.get('gps\_init\_speed\_kmh', 45.0))), 'soc0': float(sim\_cfg.get('soc0', 0.95)), 'Tb0': float(sim\_cfg.get('Tb0', 30.0)), 's0\_km': float(sim\_cfg.get('start\_s\_km', 0.0)), 'forecast\_time\_mode': str(sim\_cfg.get('forecast\_time\_mode', runtime\_cfg.get('forecast\_time\_mode', 'auto'))), 'forecast\_time\_tz': str(runtime\_cfg.get('forecast\_time\_tz', 'Australia/Darwin')), 'forecast\_start\_time\_utc': str(sim\_cfg.get('forecast\_start\_time\_utc', sim\_cfg.get('start\_utc', ''))), 'drive\_eff\_map': common\_mpc['drive\_eff\_map'], 'regen\_eff\_map': common\_mpc['regen\_eff\_map'], 'rint\_map': common\_mpc['rint\_map'], 'panel\_eff\_map': common\_mpc['panel\_eff\_map'], 'mppt\_eff\_map': common\_mpc['mppt\_eff\_map'], 'drive\_map\_eco': common\_mpc['drive\_map\_eco'], 'drive\_map\_power': common\_mpc['drive\_map\_power'], 'regen\_map\_eco': common\_mpc['regen\_map\_eco'], 'regen\_map\_power': common\_mpc['regen\_map\_power'], 'ocv\_soc\_map': common\_mpc['ocv\_soc\_map']\}]), Node(package='mpc\_solarcar', executable='mpc\_node', name='mpc\_node', parameters=[common\_mpc]), Node(package='mpc\_solarcar', executable='dashboard\_node', name='dashboard\_node', parameters=[\{'host': str(runtime\_cfg.get('dashboard\_host', '0.0.0.0')), 'port': int(runtime\_cfg.get('dashboard\_port', 8080))\}])]
              \item 主なローカル代入名: cfg, columns, common\_mpc, drive\_schedule\_yaml, forecast\_csv, key, label, path, profile\_path, profile\_yaml, required\_csvs, route\_profile\_csv, route\_waypoints\_csv, runtime\_cfg, sim\_cfg, speed\_profile\_csv, stop\_yaml
              \item 読み取る主なインスタンス属性: 特になし。
              \item 更新する主なインスタンス属性: 特になし。
              \item 明示的に送出する例外: 明示的raiseはない。
              \item 制御構造の規模: 条件分岐 0、ループ 2、try 0
            \end{itemize}
            \paragraph{この定義を読むためのPython構文}
            \begin{itemize}[leftmargin=1.6em]
            \item 特別な構文上の補足は少ない。
            \end{itemize}
            \paragraph{上から順の処理}
            \begin{enumerate}[leftmargin=1.8em]
            \item profile\_yaml に LaunchConfiguration('profile\_yaml').perform(context) の結果を代入する。
\item (profile\_path, cfg) に load\_profile(profile\_yaml) の結果を代入する。
\item runtime\_cfg に get\_section(cfg, 'runtime') の結果を代入する。
\item sim\_cfg に get\_section(cfg, 'simulation') の結果を代入する。
\item forecast\_csv に get\_path(cfg, profile\_path, 'forecast\_csv') の結果を代入する。
\item route\_waypoints\_csv に get\_path(cfg, profile\_path, 'route\_waypoints\_csv') の結果を代入する。
\item route\_profile\_csv に get\_path(cfg, profile\_path, 'route\_profile\_csv') の結果を代入する。
\item speed\_profile\_csv に get\_path(cfg, profile\_path, 'speed\_profile\_csv') の結果を代入する。
\item stop\_yaml に get\_path(cfg, profile\_path, 'stop\_yaml') の結果を代入する。
\item drive\_schedule\_yaml に get\_path(cfg, profile\_path, 'drive\_schedule\_yaml') の結果を代入する。
\item required\_csvs に \{'route waypoints': (route\_waypoints\_csv, ('dist\_km',)), 'route profile': (route\_profile\_csv, ('dist\_km',)), 'speed profile': (speed\_profile\_csv, ()), 'forecast': (forecast\_csv, ('GHI',))\} の結果を代入する。
\item ('drive\_eff\_map', 'regen\_eff\_map', 'rint\_map', 'panel\_eff\_map', 'mppt\_eff\_map', 'drive\_map\_eco', 'drive\_map\_power', 'regen\_map\_eco', 'regen\_map\_power', 'ocv\_soc\_map') を順に走査し、各要素を key に入れて処理する。
\item   required\_csvs[key] に (get\_path(cfg, profile\_path, key), ()) の結果を代入する。
\item required\_csvs.items() を順に走査し、各要素を (label, (path, columns)) に入れて処理する。
\item   require\_csv\_data\_rows(...) を実行する。
\item common\_mpc に \{'forecast\_csv': forecast\_csv, 'route\_profile\_csv': route\_profile\_csv, 'speed\_profile\_csv': speed\_profile\_csv, 'stop\_yaml': stop\_yaml, 'drive\_schedule\_yaml': drive\_schedule\_yaml, 'drive\_eff\_map': get\_path(cfg, profile\_path, 'drive\_eff\_map'), 'regen\_eff\_map': get\_path(cfg, profile\_path, 'regen\_eff\_map'), 'rint\_map': get\_path(cfg, profile\_path, 'rint\_map'), 'panel\_eff\_map': get\_path(cfg, profile\_path, 'panel\_eff\_map'), 'mppt\_eff\_map': get\_path(cfg, profile\_path, 'mppt\_eff\_map'), 'drive\_map\_eco': get\_path(cfg, profile\_path, 'drive\_map\_eco'), 'drive\_map\_power': get\_path(cfg, profile\_path, 'drive\_map\_power'), 'regen\_map\_eco': get\_path(cfg, profile\_path, 'regen\_map\_eco'), 'regen\_map\_power': get\_path(cfg, profile\_path, 'regen\_map\_power'), 'ocv\_soc\_map': get\_path(cfg, profile\_path, 'ocv\_soc\_map'), 'params\_yaml': str(profile\_path), 'profile\_runtime\_mode': 'simulation', 'forecast\_time\_mode': str(sim\_cfg.get('forecast\_time\_mode', runtime\_cfg.get('forecast\_time\_mode', 'auto'))), 'forecast\_time\_tz': str(runtime\_cfg.get('forecast\_time\_tz', 'Australia/Darwin')), 'forecast\_start\_time\_utc': str(sim\_cfg.get('forecast\_start\_time\_utc', sim\_cfg.get('start\_utc', ''))), 'soc0': float(sim\_cfg.get('soc0', 0.95)), 'Tb0': float(sim\_cfg.get('Tb0', 30.0)), 's0\_km': float(sim\_cfg.get('start\_s\_km', 0.0))\} の結果を代入する。
\item [Node(package='mpc\_solarcar', executable='gps\_sim\_node', name='gps\_sim\_node', parameters=[\{'route\_csv': route\_waypoints\_csv, 'dt': float(sim\_cfg.get('gps\_rate\_hz', 1.0)), 'init\_speed\_kmh': float(sim\_cfg.get('gps\_init\_speed\_kmh', sim\_cfg.get('v0\_kmh', 45.0)))\}]), Node(package='mpc\_solarcar', executable='solar\_state\_node', name='solar\_state\_node', parameters=[\{'profile\_yaml': str(profile\_path), 'forecast\_csv': forecast\_csv, 'route\_profile\_csv': route\_profile\_csv, 'params\_yaml': str(profile\_path), 'publish\_rate\_hz': float(sim\_cfg.get('gps\_rate\_hz', 1.0)), 'init\_speed\_kmh': float(sim\_cfg.get('v0\_kmh', sim\_cfg.get('gps\_init\_speed\_kmh', 45.0))), 'soc0': float(sim\_cfg.get('soc0', 0.95)), 'Tb0': float(sim\_cfg.get('Tb0', 30.0)), 's0\_km': float(sim\_cfg.get('start\_s\_km', 0.0)), 'forecast\_time\_mode': str(sim\_cfg.get('forecast\_time\_mode', runtime\_cfg.get('forecast\_time\_mode', 'auto'))), 'forecast\_time\_tz': str(runtime\_cfg.get('forecast\_time\_tz', 'Australia/Darwin')), 'forecast\_start\_time\_utc': str(sim\_cfg.get('forecast\_start\_time\_utc', sim\_cfg.get('start\_utc', ''))), 'drive\_eff\_map': common\_mpc['drive\_eff\_map'], 'regen\_eff\_map': common\_mpc['regen\_eff\_map'], 'rint\_map': common\_mpc['rint\_map'], 'panel\_eff\_map': common\_mpc['panel\_eff\_map'], 'mppt\_eff\_map': common\_mpc['mppt\_eff\_map'], 'drive\_map\_eco': common\_mpc['drive\_map\_eco'], 'drive\_map\_power': common\_mpc['drive\_map\_power'], 'regen\_map\_eco': common\_mpc['regen\_map\_eco'], 'regen\_map\_power': common\_mpc['regen\_map\_power'], 'ocv\_soc\_map': common\_mpc['ocv\_soc\_map']\}]), Node(package='mpc\_solarcar', executable='mpc\_node', name='mpc\_node', parameters=[common\_mpc]), Node(package='mpc\_solarcar', executable='dashboard\_node', name='dashboard\_node', parameters=[\{'host': str(runtime\_cfg.get('dashboard\_host', '0.0.0.0')), 'port': int(runtime\_cfg.get('dashboard\_port', 8080))\}])] を返す。
            \end{enumerate}
            \paragraph{代表コード断片}
            \begin{lstlisting}[language=Python]
def _setup(context):
    profile_yaml = LaunchConfiguration("profile_yaml").perform(context)
    profile_path, cfg = load_profile(profile_yaml)
    runtime_cfg = get_section(cfg, "runtime")
    sim_cfg = get_section(cfg, "simulation")

    forecast_csv = get_path(cfg, profile_path, "forecast_csv")
    route_waypoints_csv = get_path(cfg, profile_path, "route_waypoints_csv")
    route_profile_csv = get_path(cfg, profile_path, "route_profile_csv")
    speed_profile_csv = get_path(cfg, profile_path, "speed_profile_csv")
    stop_yaml = get_path(cfg, profile_path, "stop_yaml")
    drive_schedule_yaml = get_path(cfg, profile_path, "drive_schedule_yaml")

    required_csvs = {
        "route waypoints": (route_waypoints_csv, ("dist_km",)),
        "route profile": (route_profile_csv, ("dist_km",)),
        "speed profile": (speed_profile_csv, ()),
        "forecast": (forecast_csv, ("GHI",)),
    }
    for key in (
        "drive_eff_map",
        "regen_eff_map",
        "rint_map",
        "panel_eff_map",
        "mppt_eff_map",
        "drive_map_eco",
        "drive_map_power",
        "regen_map_eco",
        "regen_map_power",
        "ocv_soc_map",
    ):
        required_csvs[key] = (get_path(cfg, profile_path, key), ())
    for label, (path, columns) in required_csvs.items():
        require_csv_data_rows(path, label=label, required_columns=columns)

...
\end{lstlisting}

\subsection{L146 関数 generate\_launch\_description}
            \begin{itemize}[leftmargin=1.6em]
              \item 行範囲: L146--L152
              \item 所属: module直下
              \item 定義: \path|generate_launch_description()|
              \item docstring: 明示されていない。
              \item このブロックが直接呼ぶ主な関数・メソッド: DeclareLaunchArgument, LaunchDescription, OpaqueFunction
              \item 戻り値の要点: LaunchDescription([DeclareLaunchArgument('profile\_yaml', default\_value='config/solar/bwsc\_2027\_demo.yaml'), OpaqueFunction(function=\_setup)])
              \item 主なローカル代入名: 特になし。
              \item 読み取る主なインスタンス属性: 特になし。
              \item 更新する主なインスタンス属性: 特になし。
              \item 明示的に送出する例外: 明示的raiseはない。
              \item 制御構造の規模: 条件分岐 0、ループ 0、try 0
            \end{itemize}
            \paragraph{この定義を読むためのPython構文}
            \begin{itemize}[leftmargin=1.6em]
            \item 特別な構文上の補足は少ない。
            \end{itemize}
            \paragraph{上から順の処理}
            \begin{enumerate}[leftmargin=1.8em]
            \item LaunchDescription([DeclareLaunchArgument('profile\_yaml', default\_value='config/solar/bwsc\_2027\_demo.yaml'), OpaqueFunction(function=\_setup)]) を返す。
            \end{enumerate}
            \paragraph{代表コード断片}
            \begin{lstlisting}[language=Python]
def generate_launch_description():
    return LaunchDescription(
        [
            DeclareLaunchArgument("profile_yaml", default_value="config/solar/bwsc_2027_demo.yaml"),
            OpaqueFunction(function=_setup),
        ]
    )
\end{lstlisting}





        \subsection{launch から起動するノード}
            \begin{longtable}{p{0.07\linewidth} p{0.30\linewidth} p{0.50\linewidth}}
            \toprule
            行 & executable & package / name \\
            \midrule
            83 & \path|gps_sim_node| & package=mpc\_solarcar, name=gps\_sim\_node \\
95 & \path|solar_state_node| & package=mpc\_solarcar, name=solar\_state\_node \\
126 & \path|mpc_node| & package=mpc\_solarcar, name=mpc\_node \\
132 & \path|dashboard_node| & package=mpc\_solarcar, name=dashboard\_node \\
            \bottomrule
            \end{longtable}





        \section{処理の流れ}
        \begin{enumerate}[leftmargin=1.7em]
        \item launch 引数を受け取り、profile を読み込む。
\item Node action を動的に組み立てる。
\item LaunchDescription として ROS 2 launch へ返す。
        \end{enumerate}

        \section{このファイルの位置づけ}
        この資料群では、\path|launch/solarcar_sim.launch.py| を
        \textbf{launch}の中核として扱う。
        したがって、ここで扱う処理は単独で完結せず、
        前段の profile / launch / telemetry と後段の planner / logger / report へ繋がる。

        \end{document}
